\section{Introduction}

Channel and spatial attention blocks are among the cheapest architectural
additions in modern computer vision. A handful of extra parameters inserted
into an existing backbone is reported to buy a consistent accuracy gain, and
squeeze-and-excitation~\cite{hu2018squeeze}, the bottleneck attention
module~\cite{park2018bam}, and the convolutional block attention
module~\cite{woo2018cbam} have become canonical instances of the
pattern~\cite{guo2022attention}. Their reported gains are the reason the
applied literature routinely bolts them onto convolutional networks.

Those gains are all reported on clean test data. Each of these methods is
evaluated on top-1 and top-5 accuracy, and one of them positions itself
explicitly as an improvement over the channel-only design of another on
precisely that basis. None reports how the resulting models behave when the
input distribution shifts.

This matters because the ordering these papers induce is what a practitioner
actually consumes. Someone selecting an attention block for a deployed
classifier reads off a ranking that was measured under conditions the
deployment will not reproduce: sensors defocus, illumination drifts, contrast
collapses, images are recompressed in transit. Convolutional classifiers are
known to be brittle to exactly these
perturbations~\cite{hendrycks2019benchmarking,geirhos2019imagenet}, and a
substantial literature now treats corruption robustness as a first-class
evaluation axis~\cite{hendrycks2020augmix}. The question that follows is
therefore not whether attention blocks improve clean accuracy, but whether the
\emph{ordering} they induce on clean data is the ordering that survives
corruption.

The existing literature cannot answer this, for two reasons. First, each method
is validated in its own backbone, at its own scale, under its own training
recipe, so a cross-paper accuracy difference conflates the module with the
architecture and the optimisation schedule. Second, and more subtly, two of the
three modules considered here originate from a single shared public
implementation and inherit a common insertion convention that an independently
reimplemented third does not; even a same-codebase comparison therefore
confounds module identity with module placement.

We remove both confounds. All three attention blocks, plus an attention-free
baseline, are trained into one shared CIFAR-scale residual
backbone~\cite{he2016deep}, with each module inserted at the identical point in
the identical block, under an identical optimiser, schedule, augmentation
policy, epoch budget and seed set. Because the module is the only term that
varies, any measured difference is attributable to the module itself. We then
evaluate the trained checkpoints, without retraining, on a common-corruption
benchmark that shares the clean dataset's exact label space and
resolution~\cite{hendrycks2019benchmarking,krizhevsky2009learning}, and re-rank
the arms by mean corruption accuracy and by a relative degradation measure
rather than by clean top-1.

Our hypothesis is that the ranking may not be invariant. Channel-only
recalibration derives its gate from a globally pooled descriptor, whereas the
spatial branches of the other two modules compute attention from local feature
statistics. Spatial masks depend on where informative structure lies in the
image, and corruptions such as defocus blur and contrast reduction attack that
structure directly, so a spatial gate that mislocalises under corruption could
degrade faster than a channel-only gate. We stress in advance that the
competing outcome is equally informative: if the clean ordering is preserved at
every severity, then the field's single metric was not misleading for this
method set, which is a positive and useful finding obtainable only by running
the experiment. We commit to this symmetry before reporting results so that a
null cannot later be reframed as a failed study.

Our contributions are:

\begin{itemize}
  \item \textbf{A matched-budget, single-backbone comparison} of three
    published attention blocks against an attention-free baseline, in which
    backbone, insertion point, optimiser, schedule, augmentation and seed set
    are all held fixed, so that differences are attributable to the module
    rather than to its host or its placement.
  \item \textbf{A ranking-stability analysis} that reports whether the
    clean-accuracy ordering is preserved at each corruption severity, on both
    mean corruption accuracy and a starting-point-normalised relative
    degradation measure, with the per-corruption-type breakdown reported
    alongside the aggregate so that a skewed average is visible rather than
    hidden, and with an explicit check of whether the measured gaps exceed
    seed noise before any ordering is claimed.
  \item \textbf{A provenance record} for the compared implementations,
    including a discrepancy between the fusion rule one of these modules
    specifies in its paper and the fusion its official public code computes,
    which to our knowledge has not previously been reported.
  \item \textbf{Practical selection guidance}: whether the clean-accuracy
    ranking a practitioner would read off the original papers is the ranking
    that should drive the choice, or whether corruption robustness reorders it.
\end{itemize}

The remainder of the paper is organised as follows. Section~\ref{sec:related}
positions the work against attention-block design, corruption-robustness
benchmarking, and their recent intersection. Section~\ref{sec:method} describes
the shared backbone, the module implementations and their provenance, and the
mechanical fidelity checks. Section~\ref{sec:setup} gives the matched training
protocol and the metrics. Section~\ref{sec:results} reports the results and
answers the title question. Section~\ref{sec:limitations} states the scope of
the claim, and Section~\ref{sec:conclusion} concludes.

\section{Related Work}
\label{sec:related}

\subsection{Channel and Spatial Attention Blocks}

Squeeze-and-excitation~\cite{hu2018squeeze} introduced channel-only feature
recalibration: global average pooling reduces a feature map to a per-channel
descriptor, a two-layer bottleneck multilayer perceptron maps that descriptor
to per-channel gates, a sigmoid bounds them, and the features are rescaled
channel-wise. The design is cheap, architecture-agnostic, and has been widely
adopted; subsequent work has refined the channel-gating mechanism itself, for
example by replacing the bottleneck with a local cross-channel
interaction~\cite{wang2020eca}.

The bottleneck attention module~\cite{park2018bam} and the convolutional block
attention module~\cite{woo2018cbam} extend this by adding a spatial branch. The
former computes channel and spatial attention in parallel and fuses them; the
latter applies channel attention and then spatial attention in sequence, and
pools both average and maximum statistics in its channel gate. Both spatial
branches derive their attention maps from local feature statistics rather than
from a global descriptor, which is the property our hypothesis turns on. Broad
surveys of attention in vision confirm the pattern across the
family~\cite{guo2022attention}.

What unites this cluster for our purposes is its evaluation protocol. The
metrics reported are clean top-1 and top-5 classification accuracy, with mean
average precision added where detection is evaluated. No member of this cluster
reports a robustness or distribution-shift metric, so the relative ordering
these papers establish has never been tested away from the clean distribution.

\subsection{Corruption Robustness Benchmarking}

A parallel literature evaluates classifiers under synthetic but realistic
corruptions applied at graded severities~\cite{hendrycks2019benchmarking}. The
protocol fixes a set of corruption types, applies each at five severity levels,
and summarises degradation across the grid, which makes it possible to compare
how fast models degrade rather than only where they start. Related work
established that convolutional classifiers rely on cues that these corruptions
disrupt~\cite{geirhos2019imagenet}, and that targeted augmentation can mitigate
the damage~\cite{hendrycks2020augmix}.

The corrupted benchmark used here shares the clean dataset's exact label space
and resolution~\cite{krizhevsky2009learning}, so evaluating a trained model on
it requires inference only: no retraining, no relabelling, and no architectural
change. This is what makes the audit we perform cheap enough to be worth
running, and makes its absence from the attention-block literature harder to
explain.

Characteristically, however, this literature deploys corruption benchmarks to
\emph{validate a newly proposed intervention}. The benchmark is the evidence
that a new method helps, rather than the instrument for auditing the ranking of
methods that already exist.

\subsection{Attention Mechanisms Under Distribution Shift}

Recent work sits at the intersection of the two threads above. Attention-guided
repair procedures improve the robustness of convolutional networks against
common corruptions~\cite{ar2repair2025}, and channel attention has been
combined with multi-scale push--pull inhibition to the same
end~\cite{pushpull2025}. Both are validated on corruption benchmarks of the
kind we use.

Both, however, are proposals: each introduces a new mechanism and uses the
corruption benchmark as evidence that the new mechanism works. Neither
constitutes a head-to-head comparison of the already-published attention blocks
under a single matched training protocol with a fixed insertion point. That
controlled comparison, and the ranking-stability question it makes answerable,
is the gap this paper fills.
